% chapters/08/contoh.tex
\section{Dua Contoh Terapan}
\label{sec:contoh-bab-8}

\subsection{Contoh Pertama: Menghitung Arah dan Kecepatan secara Manual}

\begin{contohbox}
Sebuah NPC berada pada posisi $(200, 150)$ dan pemain berada pada posisi
$(500, 450)$. Kecepatan NPC ditetapkan 150 piksel per detik.

Vektor menuju pemain dihitung dari selisih kedua posisi:

\[
(500 - 200,\ 450 - 150) = (300,\ 300).
\]

Jaraknya adalah panjang vektor itu:

\[
d = \sqrt{300^2 + 300^2} = \sqrt{180000} \approx 424{,}3 \text{ piksel}.
\]

Arahnya diperoleh dengan menormalkan vektor tersebut, yaitu membagi setiap
komponen dengan panjangnya:

\[
\left(\frac{300}{424{,}3},\ \frac{300}{424{,}3}\right) \approx (0{,}707,\ 0{,}707).
\]

Kecepatannya adalah arah itu dikalikan kecepatan yang ditetapkan:

\[
(0{,}707 \times 150,\ 0{,}707 \times 150) \approx (106{,}1,\ 106{,}1)
\]

dalam satuan piksel per detik. Panjang vektor kecepatan itu kembali menjadi
150 piksel per detik, sesuai nilai yang ditetapkan.
\end{contohbox}

Apabila vektor jarak dipakai langsung sebagai kecepatan, NPC akan bergerak
dengan kecepatan 424,3 piksel per detik, yaitu hampir tiga kali lebih cepat
daripada yang ditetapkan. Itulah sebabnya penormalan diperlukan.

\subsection{Contoh Kedua: Melacak Keputusan pada Empat Jarak}

\begin{contohbox}
Sebuah NPC memiliki jari-jari persepsi 240 piksel dan jarak henti 30 piksel.
Tabel~\ref{tab:empat-jarak} memperlihatkan keputusannya pada empat jarak yang
berbeda, beserta alasan masing-masing.
\end{contohbox}

\begin{table}[htbp]
\centering
\caption{Keputusan NPC pada empat jarak terhadap pemain}
\label{tab:empat-jarak}
{\setlength{\tabcolsep}{4pt}\footnotesize%
\begin{tabularx}{\textwidth}{@{}C{2.0cm} C{2.6cm} C{3.2cm} Y@{}}
\toprule
\headrow \textbf{Jarak} & \textbf{Keputusan} & \textbf{Kecepatan} & \textbf{Alasan} \\
\midrule
400 piksel & Berdiam & 0 &
Jarak melebihi jari-jari persepsi 240 piksel, sehingga pemain tidak terlihat. \\
\addlinespace
150 piksel & Mengejar & 150 piksel/detik &
Berada di dalam jangkauan persepsi dan di luar jarak henti. \\
\addlinespace
45 piksel & Mengejar & 150 piksel/detik &
Masih di luar jarak henti, sehingga NPC tetap bergerak dengan kecepatan penuh. \\
\addlinespace
20 piksel & Berhenti & 0 &
Sudah berada di dalam jarak henti 30 piksel, sehingga NPC berhenti agar tidak
bergetar di atas pemain. \\
\bottomrule
\end{tabularx}}
\end{table}

Perhatikan baris ketiga dan keempat. Pada jarak 45 piksel NPC masih bergerak
dengan kecepatan penuh, dan pada jarak 20 piksel ia langsung berhenti. Cara itu
dipilih karena sederhana, walaupun geraknya terasa sedikit mendadak.

Apabila Anda menghendaki gerak yang lebih halus, kecepatannya dapat dikurangi
secara berangsur ketika NPC mendekati pemain, misalnya dengan mengalikan
kecepatan dengan perbandingan antara jarak dan jari-jari persepsi. Cara itu
disebut \istilah{arrival} dan menghasilkan NPC yang melambat sebelum berhenti.
Akan tetapi, cara itu menambah perhitungan dan menambah satu nilai yang harus
diatur, sehingga untuk keperluan pembelajaran pada bab ini cara yang sederhana
sudah mencukupi.

\begin{catatabox}
Keputusan NPC bergantung pada dua nilai yang Anda tetapkan,
yaitu jari-jari persepsi dan jarak henti. Keduanya tidak berasal dari
perhitungan mana pun, melainkan dari pilihan perancang. Mengubah keduanya
mengubah rasa permainan tanpa mengubah satu baris pun pada perhitungan vektor.
Itulah tanda bahwa perhitungan dan perancangan telah terpisah dengan baik.
\end{catatabox}
